% TODO(authors): confirm UAS vs ATM framing (Q-01). The motivation is written for UAS/AAM
% operations; the experiments use BlueSky transport aircraft and are described as such.
\begin{abstract}
Cooperative conflict resolution for uncrewed aircraft system (UAS) and Advanced Air Mobility
operations requires a coordination policy that keeps working when the number of vehicles $n$
in an airspace volume changes. Value-based multi-agent reinforcement learning methods that
follow centralized training with decentralized execution (CTDE), such as QMIX, use agent
and mixing networks whose input sizes are fixed by the training team size, so a trained
model cannot be evaluated at another $n$. We present GNN-QMIX, a monotonic value-decomposition
method that follows the entity-graph principle of TransfQMix and replaces its transformers
with two graph neural networks. Each vehicle builds an observation graph from the vehicles
inside its sensor radius $r_{\text{detect}}$; a shared graph convolutional network (GCN)
encodes this graph into an embedding $h_i$ and the individual action values $Q_i$. During
training only, a graph attention network (GAT) over the agent graph takes the embeddings
and the global entity state and generates the weights of a monotonic mixer, one weight row
per agent, with fixed-size mean and max readouts for the remaining terms. No parameter has
a dimension that depends on $n$, and the mixer is monotonic in every $Q_i$ for every $n$, so
each vehicle acts greedily on its own $Q_i$ at execution time and the same parameters can
be trained jointly at several fleet sizes. A second stage adds a hybrid layer in which a
PID waypoint-tracking controller issues the heading and speed commands of a vehicle and the
learned policy overrides it only while another vehicle is inside the sensor radius. We
specify an evaluation in the BlueSky air traffic simulator with joint training at
$n\in\{3,6,9\}$, zero-shot evaluation at $n\in\{10,20,50\}$, and QMIX, VDN, and IQL as
baselines. VDN and IQL are given a fixed-width observation of the three nearest aircraft so that
they, too, can be evaluated at unseen fleet sizes; QMIX, whose mixing network is tied to the
training fleet size, is trained per size. The hypothesis under test is that the proposed
method keeps its loss-of-separation performance at the unseen fleet sizes better than these
baselines.
% TODO(results): replace the last two sentences with the measured findings once
% summary_training.csv and summary_zero_shot.csv (tools/plot_results.py) exist.
\end{abstract}
